\chapter{Conclusioni}
\label{chap:conclusioni}

\section{Contributi}
\label{sec:ch7:contributi}

Questo lavoro ha costruito un sistema che rende il dato di movimento
interrogabile da un agente conversazionale senza perdere il controllo sulla sua qualità. I
contributi seguono i cinque obiettivi presentati nell'introduzione.

\paragraph{O1. Una pipeline verso un formato canonico.} Le sorgenti di
laboratorio e le riprese da telefono vengono ricondotte a diciotto centri
articolari e scritte in un C3D esteso che porta con sé provenienza, copertura,
assenze e cicli. Il file è definitivo ed il database si ricostruisce a partire da
esso, quindi i risultati sono riproducibili. Gli angoli ricalcolati riproducono
inoltre quelli di un modello clinico entro 3.2 gradi su 118 prove.

\paragraph{O2. Un archivio con la forma dei gesti.} Ogni ciclo diventa un vettore
a 32 componenti che descrive solo la forma, e su questo vettore si basano la
ricerca di gesti simili, la ripetibilità e la simmetria. La scelta dei cicli ed i
giunti esclusi per copertura insufficiente sono dichiarati. Nei casi d'uso la
ricerca dei cicli simili ha restituito per primi cinque cicli in piano della
stessa persona su 646 sessioni confrontate, mentre la ripetibilità di una ripresa
da telefono ha mostrato i limiti dei cicli ricavati dal ripiego.

\paragraph{O3. Alert dichiarativi con fonti.} Le soglie sono dati in un file
chiuso, con 76 intervalli e le loro fonti, la valutazione avviene al caricamento
ed ogni avviso riporta fonte, copertura e provenienza. Su un campione di 198
sessioni gli alert per sessione crescono con la pendenza, da 1.67 in piano a 6.02
sulla rampa al 10\%, e la stessa prova caricata come corsa passa da zero a tre
scostamenti senza che cambi nessun valore.

\paragraph{O4. Un server MCP con contratto di risposta.} Venti strumenti
restituiscono la stessa struttura, che dichiara segnale, intervallo, formula,
copertura ed avvisi, e rifiutano ciò che il dato non permette di dire. Nei tre
casi d'uso Claude e Codex hanno riportato gli stessi valori del server e gli
stessi avvisi, nessuno dei due ha stimato il segnale che mancava.

\paragraph{O5. La visualizzazione degli scostamenti.} Il manichino colora il
giunto fuori dai riferimenti solo nei fotogrammi che producono lo scostamento, e
lascia in legenda le carenze che nessun fotogramma è in grado di localizzare.

\section{Che cosa il lavoro non dimostra}
\label{sec:ch7:limiti}

Il sistema non è validato clinicamente, applica gli intervalli di riferimento in
modo tracciabile ma non ne giustifica l'uso diagnostico. Inoltre non migliora
l'accuratezza di cattura, e su quell'asse resta indietro rispetto ai sistemi
markerless. L'indipendenza dal modello linguistico riguarda i dati che l'agente
riceve, quindi il contratto garantisce che il materiale su cui l'agente lavora
contenga le proprie riserve. Che l'agente poi le riporti, e quanto vi aggiunga,
dipende invece dal modello. Nei casi d'uso entrambi gli agenti le hanno riportate,
però uno dei due ha aggiunto interpretazioni che il dato non sempre sosteneva.

\section{Prospettive}
\label{sec:ch7:prospettive}

Una volta che il software dichiara ciò che non sa, il vincolo principale diventa
la disponibilità di dati e di riferimenti. Lo sviluppo più utile è quindi quello
che allarga i dati elaborabili, cioè il riconoscimento delle parti del corpo dalla
loro geometria, seguito dall'estensione dei riferimenti ad altre condizioni ed
altri sport. La messa online con il caricamento diretto dalla chat renderebbe
infine il sistema utilizzabile da più persone, a condizione di estendere al
servizio in rete tutte le garanzie di sicurezza che oggi valgono per l'uso locale
e che sono discusse nel Paragrafo\ref{sec:ch6:mcp-online}.
